\documentclass[../../../include/open-logic-section]{subfiles}

\begin{document}
	
\olfileid{sth}{z}{nat}
\olsection{انتخاب اعداد طبیعی‌مان}

%فرض کنیم اصل غیرصوریِ~\stagesinf{} قابل توجیه است. بااین‌حال، ارزش دارد
%دربارهٔ اصل موضوع خاصی که از آن اصل غیرصوری استخراج کردیم نیز توضیحی بدهیم.

در~\olref[infinity-again]{defnomega} عبارت «اعداد طبیعی» را به‌صراحت
تعریف کردیم. این قرارداد را چگونه باید بفهمید؟ این ادعایی مابعدالطبیعی
نیست، بلکه صرفاً تصمیمی است برای آنکه مجموعه‌های معینی را اعداد طبیعی
\emph{تلقی کنیم}. در~\olref[sfr][rel][ref]{sec}،
\olref[sfr][arith][ref]{sec} و~\olref[sfr][infinite][dedekindsproof]{sec}
به دلایل این دیدگاه اشاره کردیم. اما می‌توانیم این دلایل را صریح‌تر کنیم.

اصل موضوع نامتناهی ما از~\citet{VonNeumann1925} پیروی می‌کند. اما اصل
موضوع دیگری نیز هست که می‌توانستیم به‌جای آن بپذیریم:

\begin{defish}
\emph{اصل موضوع نامتناهیِ زرملو در~\citeyear{Zermelo1908Untersuchungen}.}
مجموعه‌ای مانند~$A$ وجود دارد به‌طوری‌که~$\emptyset \in A$ و
$(\forall x \in A)\{x\} \in A$.
\end{defish}

اگر به‌جای اصل موضوع نامتناهیِ خودمان (که ملهم از فون نویمان است) از اصل
زرملو استفاده کرده بودیم، باز هم به همان اندازه یک مجموعهٔ نامتناهی
ددکیندی و در نتیجه یک جبر ددکیندی به دست می‌آوردیم. در رویکرد زرملو، عنصر
ممتاز جبر ما باز هم~$\emptyset$ (جانشین ما برای~$0$) بود، اما تزریق به‌جای
$x \mapsto x \cup \{x\}$ با نگاشت~$x \mapsto \{x\}$ داده می‌شد. ساده‌ترین
نتیجهٔ این تفاوت آن است که زرملو~$2$ را~$\{\{\emptyset\}\}$ تلقی می‌کند،
درحالی‌که ما (با فون نویمان) $2$ را
$\{\emptyset, \{\emptyset\}\}$ تلقی می‌کنیم.

چرا باید یکی از این اصول موضوع نامتناهی را به دیگری ترجیح دهیم؟ دلیل عملی
اصلی آن است که رویکرد فون نویمان برای کار با اعداد ترامتناهی به‌خوبی
«مقیاس‌پذیر» است. این موضوع را از~\olref[ordinals][]{chap} به بعد بررسی
خواهیم کرد. اما از منظر سادهٔ \emph{انجام حساب}، هر دو رویکرد به یک اندازه
کارآمدند. پس اگر کسی به شما بگوید اعداد طبیعی مجموعه \emph{هستند}، پرسش
آشکار این است: \emph{کدام مجموعه‌ها؟}

این پرسش دقیق را~\citet{Benacerraf1965} مشهور کرد. اما شایان تأکید است که
این صرفاً مشهورترین نمونهٔ پدیده‌ای است که پیش‌تر بارها با آن روبه‌رو
شده‌ایم. نکتهٔ پایه چنین است. نظریهٔ مجموعه‌ها راهی برای \emph{شبیه‌سازی}
انبوهی از انواع «شهودی» موجودات در اختیارمان می‌گذارد: البته اعداد حقیقی،
گویا، صحیح و طبیعی، اما نیز زوج‌های مرتب، تابع‌ها و رابطه‌ها. بااین‌حال،
نظریهٔ مجموعه‌ها هرگز انتخابی \emph{یکتا} از شبیه‌سازی را در اختیارمان
نمی‌گذارد. \emph{همواره} بدیل‌هایی وجود دارند که---به‌طور سرراست---به همان
اندازه به کارمان می‌آمدند.

\end{document}
